% ============================================================
%  PREÁMBULO — Informe Trabajo Integrador T.A.I.L.S.
% ============================================================
\usepackage[utf8]{inputenc}
\usepackage[T1]{fontenc}
\usepackage[spanish,es-tabla,es-noshorthands,es-nodecimaldot]{babel}
\usepackage{lmodern}
\usepackage[scaled=0.95]{helvet}
\renewcommand{\familydefault}{\sfdefault}

\usepackage[a4paper,top=2.5cm,bottom=2.8cm,left=2.5cm,right=2.5cm]{geometry}
\usepackage{graphicx}
\usepackage{float}
\usepackage{caption}
\usepackage{subcaption}
\usepackage{booktabs}
\usepackage{tabularx}
\usepackage{longtable}
\usepackage{array}
\usepackage{multirow}
\usepackage{enumitem}
\usepackage{amsmath,amssymb}
\usepackage{siunitx}
\sisetup{output-decimal-marker={,}, group-digits=false, per-mode=symbol}
\usepackage[dvipsnames,table]{xcolor}
\usepackage{listings}
\usepackage{tikz}
\usetikzlibrary{shapes.geometric,shapes.misc,arrows.meta,positioning,automata,calc,fit,backgrounds}
\usepackage{pgfplots}
\pgfplotsset{compat=1.17}
\usepackage{titlesec}
\usepackage{fancyhdr}
\usepackage{lastpage}
\usepackage[hidelinks]{hyperref}
\usepackage{url}

\setlist{itemsep=2pt,topsep=4pt}
\setlength{\parskip}{4pt}
\renewcommand{\arraystretch}{1.15}

% ---------- Encabezado / pie como la plantilla de la cátedra ----------
\pagestyle{fancy}
\fancyhf{}
\renewcommand{\headrulewidth}{0pt}
\renewcommand{\footrulewidth}{0.4pt}
\fancyfoot[C]{\footnotesize Ingeniería Mecatrónica -- Facultad de Ingeniería -- Universidad Nacional de Cuyo}
\fancyfoot[R]{\footnotesize \thepage}
\fancypagestyle{plain}{\fancyhf{}\renewcommand{\footrulewidth}{0.4pt}%
  \fancyfoot[C]{\footnotesize Ingeniería Mecatrónica -- Facultad de Ingeniería -- Universidad Nacional de Cuyo}%
  \fancyfoot[R]{\footnotesize \thepage}}

% ---------- Marcas de pendientes ----------
% Todo lo que quede en rojo debe completarse/verificarse antes de entregar.
\newcommand{\TODO}[1]{\textcolor{red}{\textbf{[PENDIENTE: #1]}}}

% ---------- Figuras con reserva de espacio ----------
% \figura[ancho]{archivo-sin-extension}{Epígrafe}{etiqueta}{Qué imagen va acá}
% Si existe figuras/<archivo>.png|.jpg|.jpeg|.pdf la inserta; si no, dibuja un
% recuadro del tamaño reservado con la indicación de qué imagen colocar.
\newcommand{\cajapendiente}[3]{%
  \fcolorbox{gray}{gray!8}{\parbox[c][#2][c]{\dimexpr#1-2\fboxsep-2\fboxrule\relax}{\centering
  \small\textcolor{gray!70!black}{\textbf{IMAGEN PENDIENTE}\\[2pt]
  \scriptsize\texttt{figuras/\detokenize{#3}}\\\texttt{(.png/.jpg/.pdf)}}}}}
\newcommand{\figura}[6][0.8\textwidth]{%
  \begin{figure}[!htbp]\centering
  \IfFileExists{figuras/#2.png}{\includegraphics[width=#1]{figuras/#2.png}}{%
  \IfFileExists{figuras/#2.jpg}{\includegraphics[width=#1]{figuras/#2.jpg}}{%
  \IfFileExists{figuras/#2.jpeg}{\includegraphics[width=#1]{figuras/#2.jpeg}}{%
  \IfFileExists{figuras/#2.pdf}{\includegraphics[width=#1]{figuras/#2.pdf}}{%
  \cajapendiente{#1}{#6}{#2}\\[2pt]{\scriptsize\textcolor{gray!70!black}{Sugerencia: #5}}}}}}
  \caption{#3}\label{#4}
  \end{figure}}
% Imagen suelta (sin float): \imagen{ancho}{alto-reservado}{archivo}
\newcommand{\imagen}[3]{%
  \IfFileExists{figuras/#3.png}{\includegraphics[width=#1]{figuras/#3.png}}{%
  \IfFileExists{figuras/#3.jpg}{\includegraphics[width=#1]{figuras/#3.jpg}}{%
  \IfFileExists{figuras/#3.jpeg}{\includegraphics[width=#1]{figuras/#3.jpeg}}{%
  \IfFileExists{figuras/#3.pdf}{\includegraphics[width=#1]{figuras/#3.pdf}}{%
  \cajapendiente{#1}{#2}{#3}}}}}}
% Versión con alto por defecto de 6 cm
\newcommand{\fig}[5][0.8\textwidth]{\figura[#1]{#2}{#3}{#4}{#5}{6cm}}

% ---------- Código ----------
\definecolor{codebg}{HTML}{F6F8FA}
\definecolor{codecom}{HTML}{6A737D}
\definecolor{codekw}{HTML}{D73A49}
\definecolor{codestr}{HTML}{032F62}
\lstdefinestyle{c}{
  language=C, backgroundcolor=\color{codebg}, basicstyle=\ttfamily\scriptsize,
  keywordstyle=\color{codekw}\bfseries, commentstyle=\color{codecom}\itshape,
  stringstyle=\color{codestr}, numbers=left, numberstyle=\tiny\color{gray},
  frame=single, rulecolor=\color{gray!40}, breaklines=true, showstringspaces=false,
  tabsize=4, captionpos=b, xleftmargin=1.5em, framexleftmargin=1.5em,
  morekeywords={uint8_t,uint16_t,uint32_t,volatile,StepperMotor,GPIO_PinState}}
\lstset{style=c}
\renewcommand{\lstlistingname}{Código}
\renewcommand{\lstlistlistingname}{Índice de códigos}

% ---------- Atajos ----------
\newcommand{\cmd}[1]{\texttt{#1}}
\newcommand{\tails}{T.A.I.L.S.}

% ---------- Estilos TikZ ----------
\tikzset{
  bloque/.style={rectangle, rounded corners=3pt, draw=black!70, fill=blue!6, thick,
                 minimum width=2.6cm, minimum height=0.9cm, align=center, font=\small},
  proceso/.style={rectangle, draw=black!70, fill=blue!5, thick, minimum width=3.4cm,
                  minimum height=0.8cm, align=center, font=\footnotesize},
  decision/.style={diamond, aspect=2.4, draw=black!70, fill=orange!12, thick, inner sep=1pt,
                   align=center, font=\footnotesize},
  terminal/.style={rounded rectangle, draw=black!70, fill=green!10, thick, minimum width=2.6cm,
                   minimum height=0.8cm, align=center, font=\footnotesize},
  flecha/.style={-{Stealth[length=2.2mm]}, thick, draw=black!70},
  estado/.style={circle, draw=black!70, thick, fill=blue!6, minimum size=2.2cm, align=center,
                 font=\footnotesize\bfseries},
}
